% !TeX root = booklet-01.tex
\section{کمیت‌های الکتریکی و خواندن اتصال‌ها}
\label{sec:foundation-circuits}
\subsection{بار، جریان و جهت مرجع}
بار الکتریکی با \(Q\) و واحد \lr{coulomb} نمایش داده می‌شود.
جریان متوسط، بار خالص عبوری از یک مقطع در واحد زمان است \cite{current}:
\begin{equation}\label{eq:foundation-current}
I_{\mathrm{avg}}=\frac{\Delta Q}{\Delta t}.
\end{equation}
برای مثال، عبور \(\qty{6}{\pico\coulomb}\) در \(\qty{2}{\nano\second}\)،
جریان متوسط \(\qty{3}{\milli\ampere}\) می‌دهد.
این رابطه از شکل تغییر جریان در طول بازه اطلاعی نمی‌دهد.
جریان لحظه‌ای با \(i(t)=\mathrm{d}Q/\mathrm{d}t\) بیان می‌شود؛ معنای نرخ تغییر
در بخش \ref{sec:foundation-math} توضیح داده شد.

جهت جریان قراردادی، جهت حرکت بار مثبت است؛ برای حرکت \lr{Electron} در فلز،
جهت حرکت حامل و جریان مخالف‌اند \cite{current}.
در تحلیل، ابتدا یک پیکان مرجع انتخاب می‌کنیم. اگر پاسخ جریان منفی شد،
جهت واقعی مخالف پیکان است؛ منفی شدن لزوماً نشانه اشتباه محاسبه نیست.

\subsection{ولتاژ، مرجع و تفاوت انرژی و توان}
ولتاژ، اختلاف پتانسیل است. تغییر انرژی پتانسیل بار مثبت \(q\) در جابه‌جایی
از نقطه \(b\) به \(a\)، برابر \(\Delta U=q(V_a-V_b)\) است.
بنابراین واحد ولتاژ، انرژی بر بار است:
\(\qty{1}{\volt}=\qty{1}{\joule\per\coulomb}\).
ولتاژ دو نقطه را با ترتیب مشخص تعریف می‌کنیم:
\begin{equation}\label{eq:foundation-voltage}
V_{ab}=V_a-V_b,\qquad V_{ba}=-V_{ab}.
\end{equation}
اگر \(V_a=\qty{1}{\volt}\) و \(V_b=\qty{0.4}{\volt}\)،
داریم \(V_{ab}=\qty{0.6}{\volt}\).
ولتاژ گره به‌تنهایی، همیشه نسبت به مرجعی مثل \lr{Ground} معنا دارد.
تغییر مرجع مشترک، اختلاف ولتاژ واقعی دو گره را تغییر نمی‌دهد.

انرژی مقدار کار قابل مبادله است؛ توان، نرخ مبادله انرژی است.
اگر جریان \(i\) وارد سر مثبت ولتاژ \(v\) یک جزء شود، توان جذب‌شده آن جزء
\(p=vi\) است. حاصل مثبت یعنی جذب و حاصل منفی یعنی تحویل انرژی به مدار.
این رابطه و کاربرد \(P=VI\) در تحلیل توان مدار با بخش \secnum{۱.۹} کتاب مرتبط‌اند
\cite{weste}. مثلاً یک مقاومت با ولتاژ \(\qty{1}{\volt}\) و جریان
\(\qty{1}{\milli\ampere}\)، توان \(\qty{1}{\milli\watt}\) جذب می‌کند؛
اگر این وضع \(\qty{2}{\milli\second}\) ثابت بماند، انرژی جذب‌شده
\(E=Pt=\qty{2}{\micro\joule}\) است. انرژی و توان هم‌واحد نیستند.

\subsection{گره، شاخه، مدار باز و اتصال کوتاه}
\lr{Node} مجموعه نقاطی است که با سیم ایدئال مستقیم متصل و هم‌پتانسیل‌اند.
وجود یک قطعه میان دو نقطه، معمولاً آن‌ها را به دو گره متفاوت تبدیل می‌کند.
\lr{Branch} یک مسیر میان دو گره است که جریان آن را در تحلیل مشخص می‌کنیم.
\lr{Ground} گره مرجع با ولتاژ صفر است؛ الزامی ندارد به زمین فیزیکی متصل باشد.

در نقشه‌های کتاب، اتصال به‌شکل \lr{T} متصل است؛ برای تقاطع دو سیم،
وجود نقطه، اتصال را مشخص می‌کند \cite[بخش \secnum{۱.۴.۲}]{weste}.
ظاهر نزدیک یا دور دو قطعه معیار هم‌گره بودن نیست؛ مسیر سیم را باید دنبال کرد.

در مدل \lr{Open Circuit}، جریان مسیر باز صفر است، ولی اختلاف ولتاژ دو سر آن
می‌تواند غیرصفر باشد. در مدل \lr{Short Circuit}، اختلاف ولتاژ صفر است؛
جریان را بقیه مدار تعیین می‌کند. اتصال کوتاه مستقیم دو سر منبع ولتاژ ایدئال
غیرصفر، مدلی ناسازگار ایجاد می‌کند؛ برای محاسبه جریان واقعی باید مقاومت‌ها و
محدودیت‌های منبع را در نظر گرفت.

\section{قوانین مدار و روش تحلیل}
\subsection{مقاومت و \lr{Ohm's Law}}
برای مقاومت خطی مثبت \(R\)، با ورود جریان مرجع به سر مثبت ولتاژ، داریم
\cite{resistors}:
\begin{equation}\label{eq:foundation-ohm}
v=Ri,\qquad i=\frac{v}{R}.
\end{equation}
مقاومت بزرگ‌تر در ولتاژ ثابت، جریان کمتری می‌دهد.
اگر جریان محاسبه‌شده منفی شود، جهت مرجع برعکس جریان واقعی بوده است.
از \(p=vi\) و رابطه \eqref{eq:foundation-ohm} نتیجه می‌گیریم
\(p=Ri^2=v^2/R\ge0\). پس مقاومت مثبت در این مدل انرژی تولید نمی‌کند.
ترانزیستور را نباید در تمام ناحیه‌های عملکرد، یک مقاومت ثابت فرض کرد؛
رابطه جریان آن در ادامه جزوه استخراج خواهد شد.

\subsection{قانون گره: \lr{KCL}}
در مدل مدار متمرکز، با حساب‌کردن جریان تمام شاخه‌ها از جمله شاخه‌های خازنی،
مجموع جریان‌های ورودی به گره با مجموع جریان‌های خروجی برابر است
\cite{kirchhoff}. این قانون با نام \lr{Kirchhoff's Current Law (KCL)} شناخته می‌شود:
\begin{equation}\label{eq:foundation-kcl}
\sum I_{\mathrm{in}}=\sum I_{\mathrm{out}}.
\end{equation}
نماد \(\sum\) یعنی جمع‌کردن همه جمله‌های مربوط.
مثلاً اگر \(\qty{3}{\milli\ampere}\) وارد گره و
\(\qty{1}{\milli\ampere}\) از یک شاخه خارج شود، جریان خروجی شاخه دیگر
\(\qty{2}{\milli\ampere}\) است. اگر خروجی به خازن متصل باشد، نباید
جریان شارژ آن را از معادله حذف کنیم.

\subsection{قانون حلقه: \lr{KVL}}
در تحلیل مدار متمرکز و با صرف‌نظر از القای مغناطیسی مدل‌نشده، جمع جبری تغییرات
ولتاژ در پیمودن یک مسیر بسته صفر است \cite{kirchhoff}.
این \lr{Kirchhoff's Voltage Law (KVL)} است:
\begin{equation}\label{eq:foundation-kvl}
\sum_{\mathrm{loop}}\Delta V=0.
\end{equation}
در عبور از منبع از پایانه منفی به مثبت، افزایش ولتاژ داریم؛ در عبور از مقاومت
در جهت جریان، افت \(IR\) داریم. مثلاً منبع \(V_s\) و دو مقاومت سری می‌دهند
\(V_s-IR_1-IR_2=0\). از اینجا
\(I=V_s/(R_1+R_2)\) به دست می‌آید. جهت پیمودن حلقه دلخواه است؛
باید علامت‌ها با همان انتخاب سازگار بمانند.

\subsection{اتصال \lr{Series} و \lr{Parallel} را از گره‌ها بشناسیم}
دو مقاومت وقتی سری‌اند که گره مشترکشان شاخه دیگری نداشته باشد؛ جریان یکسانی از
هر دو می‌گذرد. از جمع افت ولتاژها، \(R_{\mathrm{eq}}=R_1+R_2\) نتیجه می‌شود.
وقتی دو مقاومت دقیقاً میان دو گره یکسان باشند، موازی‌اند و ولتاژ یکسانی دارند.
با جمع جریان‌ها طبق \lr{KCL} به دست می‌آید \cite{resistors}:
\begin{align}
I&=\frac{V}{R_1}+\frac{V}{R_2}
=V\left(\frac{1}{R_1}+\frac{1}{R_2}\right),\notag\\
\frac{1}{R_{\mathrm{eq}}}&=\frac{1}{R_1}+\frac{1}{R_2},
\qquad R_{\mathrm{eq}}=\frac{R_1R_2}{R_1+R_2}.
\label{eq:foundation-parallel}
\end{align}
برای دو مقاومت مثبت، مقاومت معادل موازی از هر دو کمتر است؛ این یک کنترل مفید برای جواب است.
مثلاً دو مقاومت \(\qty{2}{\kilo\ohm}\) و \(\qty{3}{\kilo\ohm}\)،
در سری \(\qty{5}{\kilo\ohm}\) و در موازی \(\qty{1.2}{\kilo\ohm}\) می‌دهند.

\subsection{مثال کامل: استخراج ولتاژ یک گره}
در شکل \ref{fig:foundation-divider}، منبع ایدئال \(V_s\) و دو مقاومت مثبت داریم.
خروجی بار رسانشی ندارد. شکل برای همین جزوه رسم شده است.
\begin{figure}[ht]
\centering
\begin{latin}
\begin{circuitikz}[american]
\draw (0,0) node[ground] {} to[V,invert,l=$V_s$] (0,3)
 to[R,l=$R_1$] (3,3) to[R,l=$R_2$] (3,0) -- (0,0);
\draw (3,3) to[short,*-o] (4.7,3) node[right] {$V_x$};
\end{circuitikz}
\end{latin}
\caption{مدار آموزشی تقسیم ولتاژ؛ خروجی از گره مشترک دو مقاومت گرفته می‌شود.}
\label{fig:foundation-divider}
\end{figure}
جریان \(R_1\) را از منبع به گره \(x\) و جریان \(R_2\) را از گره به زمین
انتخاب می‌کنیم. ولتاژهای دو سر مقاومت‌ها به‌ترتیب \(V_s-V_x\) و \(V_x-0\) هستند.
با استفاده از \lr{Ohm's Law} و سپس \lr{KCL}:
\begin{align}
\frac{V_s-V_x}{R_1}&=\frac{V_x}{R_2},\notag\\
R_2(V_s-V_x)&=R_1V_x,\notag\\
R_2V_s&=(R_1+R_2)V_x,\notag\\
V_x&=V_s\frac{R_2}{R_1+R_2}.
\label{eq:foundation-divider}
\end{align}
در سطر دوم دو طرف را در \(R_1R_2\) ضرب کردیم؛ سپس پرانتز را باز و جمله‌های
حاوی \(V_x\) را یک طرف جمع کردیم. با \(V_s=\qty{1}{\volt}\)،
\(R_1=\qty{2}{\kilo\ohm}\) و \(R_2=\qty{3}{\kilo\ohm}\)،
داریم \(V_x=\qty{0.6}{\volt}\).
جریان هر دو مقاومت \(\qty{0.2}{\milli\ampere}\) و افت‌ها
\(\qty{0.4}{\volt}\) و \(\qty{0.6}{\volt}\) هستند؛ جمع افت‌ها با منبع برابر است.

\textbf{کنترل حدی و فرض بار.} با میل \(R_2\) به صفر، خروجی به صفر و با
بسیار بزرگ شدن آن نسبت به \(R_1\)، خروجی به \(V_s\) نزدیک می‌شود.
اگر بار مقاومتی \(R_L\) از خروجی به زمین وصل شود، در رابطه
\eqref{eq:foundation-divider} باید \(R_2\) را با \(R_2\parallel R_L\) جایگزین کرد؛
زیرا دیگر جریان \(R_1\) فقط از \(R_2\) نمی‌گذرد. نماد \(\parallel\) اتصال موازی را نشان می‌دهد.

\section{خازن، شارژ و مفهوم زمان}
\label{sec:foundation-capacitor}
\subsection{خازن چه چیزی را ذخیره می‌کند؟}
خازن از دو رسانا با ماده عایق میان آن‌ها ساخته می‌شود.
در مدل دوصفحه‌ای ایدئال، بارهای صفحات \(+Q\) و \(-Q\) هستند؛
مجموع بار دو صفحه می‌تواند صفر باشد، درحالی‌که انرژی در میدان ذخیره شده است
\cite{capacitors}. پس عبارت «بار خازن» معمولاً به بار علامت‌دار یک صفحه
یا اندازه بار صفحات اشاره دارد، نه مجموع آن دو.

اگر ظرفیت ثابت \(C\) باشد و \(v\) را ولتاژ صفحه اول نسبت به دوم بگیریم،
بار صفحه اول برابر است با:
\begin{equation}\label{eq:foundation-qcv}
Q=Cv,\qquad C=\frac{Q}{v}\quad(v\ne0).
\end{equation}
واحد ظرفیت \lr{farad} است و \(\mathrm{F}=\mathrm{C}/\mathrm{V}\).
مثلاً در خازن \(\qty{20}{\femto\farad}\)، ولتاژ \(\qty{1}{\volt}\)
با بار \(\qty{20}{\femto\coulomb}\) روی صفحه مثبت همراه است.
ظرفیت بزرگ‌تر در ولتاژ یکسان، بار بیشتری لازم دارد.

\subsection{از تغییر بار به تغییر ولتاژ}
برای ظرفیت ثابت، از رابطه \eqref{eq:foundation-qcv} داریم
\(\Delta Q=C\Delta v\). تقسیم دو طرف بر \(\Delta t\) و استفاده از تعریف
جریان متوسط می‌دهد:
\begin{equation}\label{eq:foundation-cap-current}
I_{\mathrm{avg}}=C\frac{\Delta v}{\Delta t},
\qquad i(t)=C\frac{\mathrm{d}v}{\mathrm{d}t}.
\end{equation}
در این قرارداد، جریان وارد صفحه‌ای می‌شود که ولتاژ مثبت برای آن تعریف شده است.
جریان مثبت بار آن صفحه و ولتاژ آن را افزایش می‌دهد؛ جریان منفی کاهش می‌دهد.
رابطه دوم شکل لحظه‌ای همان رابطه اول است. این ارتباط مبنای توضیح زمان تغییر خروجی
در کتاب و اسلاید است \cite[بخش \secnum{۱.۹}]{weste}\cite[اسلاید ۸]{mos}.
برای ظرفیت وابسته به ولتاژ، باید به رابطه عمومی \(i=\mathrm{d}Q/\mathrm{d}t\)
برگردیم؛ ثابت گرفتن \(C\) یک فرض مدل است.

\textbf{مثال آموزشی.} فرض کنیم جریان خالص ثابت \(I=\qty{10}{\micro\ampere}\)
به خازن \(C=\qty{20}{\femto\farad}\) وارد شود و ولتاژ آن از
\(\qty{0.2}{\volt}\) به \(\qty{0.8}{\volt}\) برسد. تغییر ولتاژ
\(\Delta v=\qty{0.6}{\volt}\) است. پس:
\begin{align}
\Delta t&=\frac{C\Delta v}{I}
=\frac{(20\times10^{-15}\,\mathrm{F})(0.6\,\mathrm{V})}
{10\times10^{-6}\,\mathrm{A}}\notag\\
&=1.2\times10^{-9}\,\mathrm{s}=\qty{1.2}{\nano\second}.
\label{eq:foundation-charge-time}
\end{align}
واحد \(\mathrm{F}\,\mathrm{V}/\mathrm{A}\) برابر ثانیه است.
در این مدل، دو برابر شدن ظرفیت، زمان را دو برابر و دو برابر شدن جریان، زمان را نصف می‌کند.
این نتیجه برای جریان ثابت است؛ جریان ترانزیستور واقعی عموماً طی تغییر ولتاژ ثابت نیست.

\subsection{حالت ماندگار، گذرا و پاسخ ساده \lr{RC}}
در حالت ماندگار \lr{DC}، ولتاژ خازن ثابت است؛ بنابراین در مدل ایدئال
\(\mathrm{d}v/\mathrm{d}t=0\) و جریان آن صفر می‌شود. این به معنی صفر بودن ولتاژ آن نیست.
در حالت گذرا، بار و ولتاژ تغییر می‌کنند. اختلاف ولتاژ خازن در حضور جریان محدود
نمی‌تواند ناگهان جهش کند؛ چون جهش نیازمند انتقال بار محدود در زمان صفر است.

برای درک تفاوت با جریان ثابت، خازنی را در نظر بگیرید که از منبع ثابت \(V_s\)
از طریق مقاومت \(R\) شارژ می‌شود؛ بار دیگری وجود ندارد و \(v(0)=0\).
با افزایش ولتاژ خازن، اختلاف ولتاژ مقاومت کم می‌شود؛ در نتیجه جریان
\(i=(V_s-v)/R\) کاهش می‌یابد. پس نمودار ولتاژ خط راست نیست.
از \lr{KCL} و رابطه \eqref{eq:foundation-cap-current} به دست می‌آید:
\begin{equation}\label{eq:foundation-rc-ode}
C\frac{\mathrm{d}v}{\mathrm{d}t}=\frac{V_s-v}{R}.
\end{equation}
جواب زیر فعلاً برای تفسیر معرفی می‌شود؛ می‌توان درستی آن را با مشتق‌گیری کنترل کرد:
\begin{equation}\label{eq:foundation-rc}
v(t)=V_s\left(1-e^{-t/(RC)}\right),\qquad t\ge0.
\end{equation}
عدد \(e\approx2.718\) پایه تابع نمایی است. برای عدد بی‌بعد مثبت \(x\)،
\(e^{-x}=1/e^x\)؛ با بزرگ شدن \(x\)، این مقدار به صفر نزدیک می‌شود.
حاصل \(RC\) واحد زمان دارد و آن را ثابت زمانی \(\tau=RC\) می‌نامیم.
در \(t=0\)، جواب صفر است؛ در \(t=\tau\)، مقدار ولتاژ تقریباً
\(0.632V_s\) و با زیاد شدن زمان، حد آن \(V_s\) است.
در مدل ایدئال نمایی، رسیدن دقیق به حد در زمان محدود رخ نمی‌دهد.

شکل \ref{fig:foundation-rc} همین پاسخ را بدون وابستگی به مقیاس یک فناوری
نشان می‌دهد. محور افقی زمان تقسیم بر ثابت زمانی و محور عمودی ولتاژ تقسیم بر
ولتاژ منبع است؛ هر دو نسبت بدون واحدند. کم‌شدن شیب یعنی کاهش جریان شارژ.
\begin{figure}[ht]
\centering
\begin{latin}
\begin{tikzpicture}[x=1.6cm,y=2.6cm]
\draw[->] (0,0) -- (3.6,0) node[right] {$t/\tau$};
\draw[->] (0,0) -- (0,1.25) node[above] {$v/V_s$};
\draw[dashed] (0,1) -- (3.3,1);
\node[left] at (0,1) {\rl{۱}};
\node[below left] at (0,0) {\rl{۰}};
\foreach \x/\ticklabel in {1/۱,2/۲,3/۳} {\draw (\x,0.025) -- (\x,-0.025) node[below] {\rl{\ticklabel}};}
\draw[thick,domain=0:3.3,samples=70] plot (\x,{1-exp(-\x)});
\draw[dotted] (1,0) -- (1,{1-exp(-1)}) -- (0,{1-exp(-1)});
\fill (1,{1-exp(-1)}) circle (1.4pt);
% Match the shared digit font for the Persian decimal separator.
\node[left] at (0,{1-exp(-1)}) {\fontspec[Path=C:/Windows/Fonts/]{arial.ttf}\mbox{\beginL۰٫۶۳۲\endL}};
\end{tikzpicture}
\end{latin}
\caption{نمودار تحلیلی پاسخ \lr{RC} رابطه \eqref{eq:foundation-rc}؛
این شکل از خود رابطه رسم شده است و نتیجه شبیه‌سازی \lr{SPICE} نیست.}
\label{fig:foundation-rc}
\end{figure}

برای دیدن سازگاری جواب با معادله، پاسخ و مشتق آن را کنار هم می‌نویسیم.
\[
 \begin{aligned}
 v&=V_s(1-e^{-t/\tau}),\\
 \frac{\mathrm{d}v}{\mathrm{d}t}&=\frac{V_s}{\tau}e^{-t/\tau}.
 \end{aligned}
\]
پس دو طرف رابطه \eqref{eq:foundation-rc-ode} برابر
\((V_s/R)e^{-t/\tau}\) می‌شوند. قاعده مشتق نمایی در اینجا معرفی شده است؛
استخراج تفصیلی قواعد مشتق در محل نیاز بعدی آموزش داده می‌شود.

\subsection{پرسش‌های مدار با پاسخ}
\begin{enumerate}
\item \textbf{در تقسیم ولتاژ با دو مقاومت مساوی و خروجی بی‌بار، خروجی چقدر است؟}
نصف منبع؛ زیرا \(R/(R+R)=1/2\). این حکم برای خروجی بارگذاری‌شده الزاماً برقرار نیست.
\item \textbf{اگر ولتاژ یک خازن ایدئال ثابت باشد، چه چیزی درباره جریان می‌دانیم؟}
جریان خازن صفر است، ولی شاخه‌های دیگر مدار می‌توانند جریان غیرصفر داشته باشند.
\item \textbf{چرا مدل جریان ثابت و مدل شارژ مقاومتی دو نمودار متفاوت می‌دهند؟}
در اولی شیب \(I/C\) ثابت است؛ در دومی جریان با افزایش ولتاژ خازن کاهش می‌یابد.
\item \textbf{اگر جواب جریان یک مقاومت منفی شد، چه کنیم؟}
ابتدا جهت مرجع و قطب‌های ولتاژ را کنترل کنیم؛ اگر معادله‌ها سازگار باشند،
علامت منفی فقط جهت مخالف جریان واقعی را نشان می‌دهد.
\end{enumerate}
